% Volume VII: Deep Learning Mathematics
% Continuity note for Chapter 37.
% Previous dependencies: Chapters 2-7 provide functions, matrices, affine maps, tensors, activations as elementwise maps, and matrix calculus; Chapters 31-36 provide hypothesis classes, risk, regularization, representations, and policies/value functions.
% New durable concepts: artificial neurons; affine maps plus nonlinear activations; layer shapes; batch notation; feedforward composition; parameterized function classes; ReLU and sigmoid activations; expressivity; universal approximation intuition; piecewise-linear functions; depth and width; parameter symmetries; non-identifiability; hidden-unit permutations; scaling symmetries; learned representations; feature maps; linear probes; invariance; equivariance intuition; cautious neuroscience analogies.
% Forward hooks: Chapter 38 differentiates the composed functions in this chapter; Chapter 39 studies training dynamics; Chapter 40 specializes architectures; Chapter 42 analyzes representation geometry, robustness, and interpretability; Chapter 48 compares artificial and neural population representations.
% Notation commitments: Inputs are x in R^{d_0}; layer preactivations are z_\ell; activations are h_\ell in R^{d_\ell}; h_0=x; weights are W_\ell in R^{d_\ell x d_{\ell-1}}; biases are b_\ell in R^{d_\ell}; elementwise nonlinearities are phi_\ell; full network is f_\theta.
\section{Chapter 37: Neural Networks as Parameterized Function Classes}
\label{ch:neural-networks-parameterized-function-classes}

The previous machine-learning chapters treated a model class \(\mathcal F\) abstractly. A neural network is one important way to build such a class. It is a composition of simple parameterized maps: affine transformations, elementwise nonlinearities, normalizations, pooling operations, attention mechanisms, or other differentiable modules. This chapter begins with the simplest feedforward version because the same shape discipline extends to later architectures.

The central question is not whether neural networks are inspired by the brain in a loose sense. The mathematical question is what functions they can represent, how parameters determine those functions, why different parameter settings can compute the same function, and how hidden layers create learned feature maps.

\paragraph{Chapter problem.}
How do artificial neurons and layers define parameterized function classes, what makes those classes expressive, and how should we interpret hidden activations as representations?

\paragraph{Section map.}
Section 37.1 defines artificial neurons, affine layers, activations, batches, and network composition. Section 37.2 studies expressivity through nonlinear activations, piecewise-linear examples, and universal approximation intuition. Section 37.3 explains parameter symmetries and non-identifiability. Section 37.4 treats hidden layers as learned representations for downstream prediction, invariance, and neural population analogies.

\subsection{37.1 Artificial neurons and layers}

\paragraph{Guiding question.}
What mathematical map is computed by one artificial neuron, one layer, and a feedforward network?

\paragraph{Beginner-facing intuition.}
An artificial neuron takes a weighted sum of inputs, adds a bias, and passes the result through a nonlinearity. A layer repeats this operation for many output units at once. A network composes layers so that later features are functions of earlier features.

The linear algebra is ordinary. The important new habit is shape discipline. Each layer has an input dimension, output dimension, weight matrix, bias vector, and activation function. If those shapes do not match, the network is not a well-defined function.

\paragraph{Formal definitions.}
Let the input be
\[
x\in\mathbb R^{d_0}.
\]
A single artificial neuron has weights \(w\in\mathbb R^{d_0}\), bias \(b\in\mathbb R\), activation \(\phi:\mathbb R\to\mathbb R\), and output
\[
\boxed{
h=\phi(w^\top x+b).
}
\]
A layer with \(d_1\) output units has
\[
W_1\in\mathbb R^{d_1\times d_0},\qquad b_1\in\mathbb R^{d_1},
\]
preactivation
\[
z_1=W_1x+b_1\in\mathbb R^{d_1},
\]
and activation
\[
h_1=\phi_1(z_1)\in\mathbb R^{d_1},
\]
where \(\phi_1\) is applied elementwise unless stated otherwise.

A feedforward network with \(L\) layers is
\[
h_0=x,
\qquad
z_\ell=W_\ell h_{\ell-1}+b_\ell,
\qquad
h_\ell=\phi_\ell(z_\ell),
\]
for \(\ell=1,\ldots,L\). The network output is
\[
f_\theta(x)=h_L
\]
or sometimes an affine readout from \(h_L\). The parameter vector \(\theta\) collects all weights and biases.

For a batch \(X\in\mathbb R^{B\times d_0}\), with rows as examples, a layer is often implemented as
\[
Z=XW^\top+\mathbf 1 b^\top\in\mathbb R^{B\times d_1}.
\]

\paragraph{Core mechanism: layers compose typed functions.}
Each layer is a function
\[
F_\ell:\mathbb R^{d_{\ell-1}}\to\mathbb R^{d_\ell},
\qquad
F_\ell(h)=\phi_\ell(W_\ell h+b_\ell).
\]
The whole feedforward network is the composition
\[
\boxed{
f_\theta=F_L\circ F_{L-1}\circ\cdots\circ F_1.
}
\]
The codomain of \(F_{\ell-1}\) must equal the domain of \(F_\ell\). Thus \(W_\ell\) must have \(d_{\ell-1}\) columns and \(d_\ell\) rows. This typed composition is what later backpropagation differentiates.

\paragraph{Concrete example.}
Let
\[
x=\begin{pmatrix}2\\-1\end{pmatrix},\qquad
w=\begin{pmatrix}3\\1\end{pmatrix},\qquad
b=-2,
\]
and use ReLU
\[
\phi(u)=\max(0,u).
\]
Then
\[
w^\top x+b=3(2)+1(-1)-2=3,
\]
so the neuron output is
\[
h=\max(0,3)=3.
\]
If the input were \(x=(0,1)^\top\), the preactivation would be \(-1\), and the ReLU output would be \(0\).

\paragraph{Computational neuroscience example.}
A firing-rate model may write a neuron's mean response as
\[
r=\phi(w^\top x+b),
\]
where \(x\) is a stimulus feature vector and \(\phi\) enforces nonnegative firing rates. This resembles an artificial neuron mathematically, but it is a simplified input-output model, not a full biophysical model of spikes, dendrites, conductances, or recurrent circuitry.

\paragraph{AI and machine-learning example.}
In image classification, a network maps an image tensor to class scores. Early layers may compute local features, later layers combine them, and the final layer produces logits in \(\mathbb R^K\). The softmax converts logits to class probabilities. Every module still has a typed input and output shape.

\paragraph{Common misconceptions and failure modes.}
A layer is not just a matrix multiplication; without nonlinearities, a stack of affine layers collapses to one affine map. Biases matter because they shift activation thresholds. Batch dimensions are not feature dimensions. An artificial neuron is a mathematical unit in a model, not a literal biological neuron. Shape errors often arise from transposing \(W\) inconsistently between row-batch and column-vector conventions.

\paragraph{Exercises.}
\begin{enumerate}[label=\textbf{37.1.\arabic*.}, leftmargin=*]
\item \textbf{Mechanical.} If \(h_{\ell-1}\in\mathbb R^{64}\) and \(h_\ell\in\mathbb R^{32}\), state the shapes of \(W_\ell\) and \(b_\ell\) in column-vector convention.
\item \textbf{Proof or derivation.} Show that if all activations are identity maps, then a two-layer affine network is equivalent to a single affine map.
\item \textbf{Numerical/computational.} Compute \(\phi(w^\top x+b)\) for \(x=(1,2)^\top\), \(w=(2,-1)^\top\), \(b=0.5\), and ReLU activation.
\item \textbf{Interpretation.} In a firing-rate model \(r=\phi(w^\top x+b)\), what do \(x\), \(w\), \(b\), and \(r\) represent?
\item \textbf{Misconception/debugging.} A student uses \(W\in\mathbb R^{d_{\mathrm{in}}\times d_{\mathrm{out}}}\) in the formula \(z=Wh+b\) with \(h\in\mathbb R^{d_{\mathrm{in}}}\). Identify the shape problem and one possible convention that fixes it.
\end{enumerate}

\subsection{37.2 Expressivity}

\paragraph{Guiding question.}
What kinds of functions can a neural network represent, and why do nonlinear activations and depth matter?

\paragraph{Beginner-facing intuition.}
Expressivity is about the functions available to a model class before training. A linear model can only draw linear boundaries or fit affine trends. A neural network with nonlinear activations can bend, threshold, combine, and reuse features. More units or layers enlarge the set of possible functions, but expressivity alone does not guarantee learnability or generalization.

The simplest way to see expressivity is one-dimensional ReLU networks. A ReLU is flat on one side and linear on the other. Sums of shifted ReLUs create piecewise-linear curves with bends at chosen thresholds.

\paragraph{Formal definitions.}
A parameterized network class is
\[
\mathcal F=\{f_\theta:\mathcal X\to\mathcal Y,\ \theta\in\Theta\}.
\]
Its expressivity concerns which target functions can be represented exactly or approximated well by members of \(\mathcal F\). Approximation is usually measured by a norm, such as
\[
\|f-g\|_\infty=\sup_{x\in K}|f(x)-g(x)|
\]
on a compact set \(K\), or an \(L^2(P)\) error under a data distribution.

For ReLU
\[
\sigma(t)=\max(0,t),
\]
a one-hidden-layer scalar-output network has form
\[
f(x)=\sum_{j=1}^m a_j\sigma(w_jx+b_j)+c.
\]
In one dimension, this is a continuous piecewise-linear function.

\paragraph{Core theorem: universal approximation, stated carefully.}
Let \(K\subset\mathbb R^d\) be compact and let \(\phi\) be a suitable non-polynomial activation such as sigmoid or ReLU. Then, for every continuous function \(g:K\to\mathbb R\) and every \(\varepsilon>0\), there exists a one-hidden-layer network \(f_\theta\) with sufficiently many hidden units such that
\[
\sup_{x\in K}|f_\theta(x)-g(x)|<\varepsilon.
\]
This theorem says that the class is dense in continuous functions on compact sets. It does not say that a small network suffices, that gradient descent will find the approximation, that finite data identify it, or that the approximation generalizes.

\paragraph{Worked mechanism: ReLUs build bends.}
On the real line,
\[
f(x)=\sigma(x)-2\sigma(x-1)+\sigma(x-2)
\]
is
\[
f(x)=
\begin{cases}
0, & x<0,\\
x, & 0\leq x<1,\\
2-x, & 1\leq x<2,\\
0, & x\geq2.
\end{cases}
\]
The first ReLU creates a positive slope at \(0\), the second changes the slope by \(-2\) at \(1\), and the third changes the slope by \(+1\) at \(2\). The result is a triangular bump. Adding more shifted ReLUs creates more breakpoints.

\paragraph{Concrete example.}
The absolute-value function can be represented exactly by two ReLUs:
\[
|x|=\sigma(x)+\sigma(-x).
\]
For \(x=3\), the value is \(3+0=3\). For \(x=-2\), the value is \(0+2=2\). A linear model \(ax+b\) cannot equal \(|x|\) on all of \([-1,1]\) because \(|x|\) has different slopes on the two sides of zero.

\paragraph{Computational neuroscience example.}
Neural response functions can be nonlinear in stimulus features. A network can approximate nonlinear tuning curves or decision boundaries from population activity. The expressivity theorem only says such functions can be approximated mathematically; it does not imply that biological circuits use the same parameterization or training algorithm.

\paragraph{AI and machine-learning example.}
Deep networks can represent hierarchical features: edges combine into parts, parts into objects; tokens combine into phrases, phrases into longer contexts. Depth can sometimes represent compositional functions with far fewer units than a shallow network, though exact depth-separation results require assumptions about the target function class.

\paragraph{Common misconceptions and failure modes.}
Universal approximation is not universal learning. It ignores optimization, sample size, noise, computation, and inductive bias. More expressive models can overfit or learn shortcuts. A network with no nonlinearity is still linear after composition. A function being representable does not mean its parameters are unique.

\paragraph{Exercises.}
\begin{enumerate}[label=\textbf{37.2.\arabic*.}, leftmargin=*]
\item \textbf{Mechanical.} Write a one-hidden-layer ReLU network with two hidden units that represents \(|x|\).
\item \textbf{Proof or derivation.} Verify piece by piece that \(\sigma(x)-2\sigma(x-1)+\sigma(x-2)\) is the triangular bump described above.
\item \textbf{Numerical/computational.} Evaluate \(f(x)=\sigma(x)+\sigma(-x)\) at \(x=-3,0,4\).
\item \textbf{Interpretation.} Why does a universal approximation theorem not guarantee good performance on a small training set?
\item \textbf{Misconception/debugging.} A student says, ``A deep linear network is expressive because it has many layers.'' Show why stacked linear maps still produce a linear map.
\end{enumerate}

\subsection{37.3 Parameter symmetries and identifiability}

\paragraph{Guiding question.}
Can different neural-network parameter values compute exactly the same input-output function?

\paragraph{Beginner-facing intuition.}
In many statistical models, different parameters imply different distributions or functions. Neural networks often violate this property. Hidden units can be permuted. ReLU units can be rescaled across adjacent layers. Overparameterized networks can contain flat or nearly flat directions in parameter space. Thus a learned parameter vector is not always identifiable from the function it computes.

This matters for optimization and interpretation. Looking at one weight in isolation may be misleading if many equivalent parameterizations exist.

\paragraph{Formal definitions.}
Let
\[
\theta\mapsto f_\theta
\]
be the map from parameters to functions. Parameters are identifiable from functions if
\[
f_\theta=f_{\theta'} \quad \Longrightarrow \quad \theta=\theta'
\]
or equal up to a declared equivalence. A parameter symmetry is a transformation \(T:\Theta\to\Theta\) such that
\[
\boxed{
f_{T(\theta)}=f_\theta
}
\]
for all relevant inputs, even though \(T(\theta)\neq\theta\).

\paragraph{Core mechanism: hidden-unit permutation symmetry.}
Consider a one-hidden-layer network
\[
f(x)=\sum_{j=1}^m a_j\phi(w_j^\top x+b_j)+c.
\]
Let \(\pi\) be any permutation of \(\{1,\ldots,m\}\). Define new parameters
\[
a'_j=a_{\pi(j)},\qquad w'_j=w_{\pi(j)},\qquad b'_j=b_{\pi(j)}.
\]
Then
\[
\sum_{j=1}^m a'_j\phi((w'_j)^\top x+b'_j)
=\sum_{j=1}^m a_{\pi(j)}\phi(w_{\pi(j)}^\top x+b_{\pi(j)})
=\sum_{k=1}^m a_k\phi(w_k^\top x+b_k).
\]
The function is unchanged because addition is commutative. Therefore the hidden units have no intrinsic ordering.

\paragraph{Scaling symmetry for ReLU.}
ReLU is positively homogeneous:
\[
\sigma(\alpha u)=\alpha\sigma(u)\qquad \text{for }\alpha>0.
\]
Thus
\[
a_j\sigma(w_j^\top x+b_j)
=\frac{a_j}{\alpha}\sigma((\alpha w_j)^\top x+\alpha b_j)
\]
for any \(\alpha>0\). Scaling incoming weights and bias by \(\alpha\) while scaling the outgoing weight by \(1/\alpha\) leaves the function unchanged.

\paragraph{Concrete example.}
Let
\[
f(x)=2\sigma(x-1)+3\sigma(-x).
\]
Swapping the two hidden units gives
\[
f(x)=3\sigma(-x)+2\sigma(x-1),
\]
the same function. Also,
\[
2\sigma(x-1)=1\cdot\sigma(2x-2),
\]
so the first hidden unit has a different parameterization with the same contribution.

\paragraph{Computational neuroscience example.}
Latent neural-network models fitted to neural data may contain hidden units or factors whose ordering, signs, or scales are arbitrary. A rotated or permuted latent representation can fit the same data equally well. Interpreting a single hidden coordinate as a specific biological variable requires constraints, validation, or alignment with independent measurements.

\paragraph{AI and machine-learning example.}
Permutation symmetry creates many equivalent optima in neural-network loss landscapes. Scaling symmetries affect the size of weights and gradients. Normalization, weight decay, and parameterization choices can break or soften some symmetries, but they do not remove all non-identifiability.

\paragraph{Common misconceptions and failure modes.}
Different parameters do not necessarily mean different learned functions. A large weight norm may be partly a scaling artifact. Hidden-unit labels are arbitrary unless a convention is imposed. Non-identifiability is not always a bug; it can coexist with a well-defined predictive function. But it makes mechanistic interpretation of parameters harder.

\paragraph{Exercises.}
\begin{enumerate}[label=\textbf{37.3.\arabic*.}, leftmargin=*]
\item \textbf{Mechanical.} For a two-hidden-unit network, write the parameters after swapping hidden units \(1\) and \(2\).
\item \textbf{Proof or derivation.} Prove the ReLU scaling symmetry \(a\sigma(w^\top x+b)=(a/\alpha)\sigma((\alpha w)^\top x+\alpha b)\) for \(\alpha>0\).
\item \textbf{Numerical/computational.} Evaluate \(2\sigma(x-1)\) and \(\sigma(2x-2)\) at \(x=0,1,3\) and compare.
\item \textbf{Interpretation.} Why can permutation symmetry make it difficult to compare hidden units across two independently trained networks?
\item \textbf{Misconception/debugging.} A student says, ``If two trained networks have different weights, they must have learned different functions.'' Give a symmetry-based correction.
\end{enumerate}

\subsection{37.4 Neural networks as learned representations}

\paragraph{Guiding question.}
How do hidden layers turn raw inputs into feature spaces where downstream tasks become easier?

\paragraph{Beginner-facing intuition.}
A hidden layer is not only an intermediate computation. It is a representation map. It sends an input \(x\) to a feature vector \(h_\ell(x)\). A good representation may separate classes, preserve task-relevant variation, discard nuisance variation, or make dynamics simpler. The same network can be read as a predictor and as a sequence of learned coordinate systems.

Representation learning is powerful because the final classifier or decoder may be simple even when the raw input-output relation is complicated. The difficulty is that ``good representation'' depends on the downstream task and invariances we want.

\paragraph{Formal definitions.}
For a network
\[
h_0(x)=x,\qquad h_\ell(x)=\phi_\ell(W_\ell h_{\ell-1}(x)+b_\ell),
\]
the map
\[
\boxed{
h_\ell:\mathcal X\to\mathbb R^{d_\ell}
}
\]
is the layer-\(\ell\) representation. A linear probe fits a simple readout
\[
g(h)=Ah+c
\]
or a softmax classifier on frozen representations \(h_\ell(x_i)\). If a task is linearly decodable from \(h_\ell\), then relevant information is accessible to a linear readout at that layer.

An invariant representation satisfies approximately
\[
h(Tx)\approx h(x)
\]
for transformations \(T\) that should not change the task label, such as small image translations in some classification tasks. An equivariant representation changes predictably:
\[
h(Tx)\approx \widetilde T h(x).
\]

\paragraph{Core mechanism: nonlinear features can make linear readouts sufficient.}
The XOR labels on inputs \((0,0),(1,0),(0,1),(1,1)\) are not linearly separable in raw coordinates. Define a feature map
\[
h(x_1,x_2)=
\begin{pmatrix}
x_1\\
x_2\\
x_1x_2
\end{pmatrix}.
\]
The classifier
\[
g(h)=x_1+x_2-2x_1x_2
\]
equals \(1\) on \((1,0)\) and \((0,1)\), and \(0\) on \((0,0)\) and \((1,1)\). Thus a nonlinear representation can turn a nonlinear raw task into a linear readout problem. Neural networks learn such feature maps from data rather than requiring them to be manually specified.

\paragraph{Concrete example.}
For \(x=(1,0)\),
\[
h(x)=(1,0,0)^\top,\qquad g(h)=1.
\]
For \(x=(1,1)\),
\[
h(x)=(1,1,1)^\top,\qquad g(h)=1+1-2=0.
\]
The product feature \(x_1x_2\) supplies information that a raw linear classifier lacked.

\paragraph{Computational neuroscience example.}
Population activity can be viewed as a representation of stimuli, actions, or latent task variables. A linear decoder asks whether a variable is linearly readable from neural activity. This is analogous to a linear probe on artificial representations. The analogy is useful but limited: neural activity is produced by a biological circuit with recurrent dynamics, noise, and behavioral context, not by a supervised deep network alone.

\paragraph{AI and machine-learning example.}
In transfer learning, an encoder trained on large data is frozen or fine-tuned for new tasks. If the representation clusters semantically similar inputs and makes labels linearly separable, a small downstream dataset may be enough. In reinforcement learning, a learned state representation can make value functions and policies easier to approximate.

\paragraph{Common misconceptions and failure modes.}
Linear decodability does not prove that a biological system uses the decoded variable. Invariance can discard information that another task needs. A representation may encode spurious shortcuts such as background or recording artifacts. High-dimensional representations can make many variables decodable by chance without proper controls. Similar-looking embeddings can hide different causal mechanisms.

\paragraph{Exercises.}
\begin{enumerate}[label=\textbf{37.4.\arabic*.}, leftmargin=*]
\item \textbf{Mechanical.} For a layer \(h_\ell:\mathbb R^{100}\to\mathbb R^{20}\), state the shape of a linear readout \(A\) that predicts \(5\) logits from \(h_\ell(x)\).
\item \textbf{Proof or derivation.} Verify that \(g=x_1+x_2-2x_1x_2\) implements XOR on binary inputs when the positive class is exactly one active bit.
\item \textbf{Numerical/computational.} Compute \(h(x)\) and \(g(h)\) for all four binary inputs in the XOR example.
\item \textbf{Interpretation.} What does it mean, and not mean, if reach direction is linearly decodable from a neural population vector?
\item \textbf{Misconception/debugging.} A student says, ``If a network representation is invariant to a transformation, that is always good.'' Give a task where that invariance would hurt.
\end{enumerate}

Neural networks are therefore not mysterious black boxes at the mathematical level. They are typed compositions of affine maps and nonlinear operations, forming expressive but non-identifiable function classes whose hidden states can be used as representations. The next chapter differentiates these compositions efficiently, turning the functions of this chapter into trainable differentiable programs.

\clearpage
